% !TeX root = ../papers/nc-p3.tex
\section{倍率反転と階数15の座標端の一般同型類}
\label{sec:rank-fifteen-endpoint-orbits}
\begin{theorem}\label{thm:rank-fifteen-endpoint-orbits}
標数零の通常な階数15窓を考える。第一表示を $(z,\lambda,n)$ で取ると、
第二表示の $(z,1/(z\lambda),n)$ と添字保持同型である。
両表示が倍率を動かして与える添字保持同型類の集合は一致する。
また第一表示内で $\sigma(z,\lambda,n)=(1/z,1/\lambda,n)$ は添字保持同型を与える。
\end{theorem}
\begin{proof}
$k,m$ は $\lambda\mapsto1/(z\lambda)$ で不変であり、$\sigma$ ではともに $z$ 倍となる。
定理\ref{thm:rank-fifteen-covers}の $P(n)$ は両変換で同一である。
前者では第一表示の枠は同じで、対象の第二表示の枠は $(zb,a,zf,e)$ となる。
後者では対象の第一表示の枠が $(a,zb,e,zf)$ となる。
$r^2=z,d^2=\lambda$ とする。前者の対象テンソルは、元のテンソルの
$(i,j,k\oplus1,t\oplus1)$ 成分に、各位置で次の対角係数を掛けて得られる。
\[
 (rd,rd,1,1),\quad(rd,rd,1,1),\quad
 ((rd)^{-1},(rd)^{-1},z^{-1},1),\quad(r/d,(rd)^{-1},1,1).
\]
後者では四文字に同じ置換 $(2\;3)$ を施し、位置ごとに
\[
 (d/r,d/r,z^{-1},z^{-1}),\quad(rd,d/r,1,1),\quad
 (r/d,r/d,1,1),\quad((rd)^{-1},r/d,1,1)
\]
を掛ける。\path{computations/verify_rank_fifteen_word_symmetries.py} は
枠の四成分を独立変数として両恒等式を検算する。四位置で繰り返す矢の基底変更が
すべての隣接関係を対応させる。いずれも位置の順序と始点を保つ。
\end{proof}

\begin{theorem}\label{thm:endpoint-generic-fibers}
前節の第一表示の楕円曲線 $n=1$ 内で、一般点と添字保持同型なパラメータは
元の点とその $\sigma$ 像のちょうど二点である。
有理曲線 $m=ik,n=i$ の各成分内でも同様で、その二点は $t,-1/t$ である。
楕円曲線と有理曲線の添字保持同型類の像は異なる曲線である。
\end{theorem}
\begin{proof}
$T$ を節\ref{sec:rank-fifteen-adjoint-invariants}の写像とし、スカラー行列を零に送って
全行列空間に拡張する。枠の規格化 $a+e=2$ において反対切断の左右の核行列は
それぞれ $\operatorname{diag}(1,1,-1,-1/z)$ と $\operatorname{diag}(1,1,-1,-1)$ であり、
切断に掛けて確認できる。
$T$ をトレース形式で $H\in\operatorname{End}(V_0)\otimes\operatorname{End}(V_1)$
と見る。$P=(4,0,3,1,2)$、$Q=(2,4,3,0,1)$ を $0,\ldots,4$ での像の列とし、
\[
 D=\sum_{a,b\in\{0,1,2,3\}^5}
       \prod_{j=0}^4 H_{a_j,a_{P(j)},b_j,b_{Q(j)}},\qquad
 L=\frac{D^2}{(\operatorname{tr} T^*T)^5}
\]
と置く。各行列因子の行・列は置換の閉路に沿って縮約されるので、二つの独立共役を保ち、
$T$ のスカラー倍も $L$ では相殺する。従って $L$ は添字保持不変量である。

\path{computations/verify_rank_fifteen_endpoint_orbits.py} は各曲線の関数体上で
$J_2,J_3$ を計算する。楕円曲線では二式とも $t$ だけの関数になる。$t=2$ の不変量と
一致するための二式の分子のモニック最大公約式は、楕円・有理の順に
\[
 (t-2)(t-1/2)(t^2-110t/9+1),\qquad
 (t-2)(t-3)(t+1/2)(t+1/3).
\]
無限遠は最高次係数から除かれる。楕円上の $t=2$ の値を有理曲線上で探すと
最大公約式は1で無限遠にも解がなく、両像の相違が従う。

楕円上の $t=2,Y^2=46$ で $Y\equiv47\pmod{103}$ の剰余を取る。
上の四根は103で $2,52,25,33$ となり、それぞれ二つの $Y$ も相異なる。
$L$ は初めの二根上で99、残りの二根上で71となる。
命題\ref{prop:rank-fifteen-oriented-cubic}の $C$ は初めの二根の各々で65と63になり、
各一方だけが指定点と一致する。有理曲線では $i\equiv10\pmod{101}$ とし、
$t=2,-1/2$ で $C=93$、$t=3,-1/3$ で $C=98$ を得る。共役な $i$ の成分にも
同じ議論を適用できる。全分母は選んだ剰余点で単元である。

両最大公約式は平方自由であり、楕円の二次被覆もこの各点で不分岐である。
滑らかな完備曲線から有理不変量の値の射影直線への写像を取れば、
このファイバーは既に $J_2,J_3$ の段階で被約である。
$C,L$ を加えて残る長さは特殊点で高々2であるから、有限写像の一般ファイバーの
長さも高々2である。前定理の $\sigma$ は一般点を動かし、実際に二点を同型にする。
この比較は各正規化成分内のもので、有理曲線の二つの符号成分間の比較は残る。
\end{proof}
